% Budget: ~1 page. Source of truth: src/problem.py (objectives, split),
% src/build_instances.py and src/build_cvrplib_instances.py (data generation).
% Every constant below is taken from those files; update both together.
\section{Many-Objective Waste-Collection Routing Model}
\label{sec:problem}

\subsection{Setting}
We model municipal waste collection as a capacitated vehicle routing problem
(CVRP)~\cite{toth2014vrp,belien2014waste}. A directed graph $G=(V,A)$ has
node $0$ as the depot and customers $C=V\setminus\{0\}$, $|C|=n$, each a
collection point with waste demand $q_i>0$. A homogeneous fleet of trucks with
capacity $Q$ starts and ends at the depot. The fleet size is not fixed: the
number of routes $K$ is a consequence of the solution and is itself traded off
through the cost objective. Each arc $(i,j)\in A$ carries a distance $d_{ij}$
(km); the Sibiu matrices are asymmetric, as distances on a real directed
road network generally are.

\subsection{Solution representation and decoding}
Every algorithm in this study searches the same space: a \emph{giant tour}
$\pi$, i.e., a permutation of $C$. A deterministic capacity-first greedy split
walks $\pi$ in order and closes the current route whenever adding the next
customer would exceed $Q$, in the spirit of route-first/cluster-second
decoding~\cite{prins2004simple}. Route $k$ is then the node sequence
$r_k=(0,v_{k,1},\dots,v_{k,m_k},0)$. Because the split never lets a route's
load exceed $Q$, every permutation decodes to a feasible solution: no repair
operator and no penalty term are needed anywhere in the pipeline. This matters
for the qubit encoding (\cref{sec:algorithm}), where an earlier single-objective
QIEA on the same Sibiu routes lost population diversity on large instances
because of its repair step~\cite{gilea2026scalable}.

\subsection{Objectives}
Let $A(r_k)$ be the arcs traversed by route $k$ and $\mathcal{R}=\{r_1,\dots,r_K\}$
the decoded routes. Travel time and variable cost are derived from distance
through a per-arc \emph{road factor} $\rho_{ij}>0$ that models congestion and
road type independently of length:
\begin{equation}
t_{ij}=\frac{\rho_{ij}\,d_{ij}}{\bar v},\qquad c_{ij}=c_v\,\rho_{ij}\,d_{ij},
\end{equation}
with base urban speed $\bar v=25$~km/h and $c_v=1$ cost unit per km. All five
objectives are minimized:
\begin{align}
f_1 &= \sum_{k=1}^{K}\sum_{(i,j)\in A(r_k)} d_{ij}
  && \text{(distance, km)}\\
f_2 &= \sum_{k=1}^{K} T_k,\quad T_k=\sum_{(i,j)\in A(r_k)} t_{ij}
  && \text{(time, h)}\\
f_3 &= c_F\,K+\sum_{k=1}^{K}\sum_{(i,j)\in A(r_k)} c_{ij}
  && \text{(cost)}\\
f_4 &= \sum_{k=1}^{K}\sum_{(i,j)\in A(r_k)}
  \bigl(e_d\,d_{ij}+e_t\,t_{ij}+e_q\,d_{ij}\,L_{k,j}\bigr)
  && \text{(CO$_2$, kg)}\\
f_5 &= \sqrt{\tfrac{1}{K}\textstyle\sum_{k=1}^{K}\bigl(T_k-\bar T\bigr)^2}
  && \text{(workload balance)}
\end{align}
where $c_F=40$ is the fixed cost per vehicle used, $\bar T$ is the mean route
time, and $f_5=0$ when $K=1$. The emission model has a distance term
($e_d=0.7$~kg/km), an idling term that grows with congestion
($e_t=1.2$~kg/h), and a payload term ($e_q=0.004$~kg per km per kg of load).
$L_{k,j}$ is the truck's cumulative load on route $k$ after collecting at
node $j$ (the full route load on the final arc back to the depot). Because a
collection truck fills up as it drives, $f_4$ depends on the \emph{order} of
visits, not only on which arcs are used. $f_3$ couples to fleet size through
$c_F$, and $f_5$ rewards splitting work evenly across trucks, which conflicts
with using few, long routes.

The road factor is essential rather than cosmetic. In an early version of the
model, time, cost and emissions were scalar multiples of distance; the five
objectives were then perfectly correlated and every Pareto front collapsed to a
single point. Drawing $\rho_{ij}$ independently per arc,
$\rho_{ij}\sim\mathrm{LogNormal}(0,\,0.35^2)$, removes this collinearity.

\subsection{Instances}
\Cref{tab:instances} lists the seven instances. The three Sibiu instances use
real road-network distance matrices for waste-collection routes in Sibiu,
Romania, with 333, 198 and 201 collection points. No municipal demand records
are available for these routes, so demands are synthesized as
$q_i\sim\mathcal{U}\{40,\dots,400\}$~kg and capacity is set to
$Q=\lceil 1.15\sum_i q_i/6\rceil$, i.e., about six trucks with 15\% slack. To
test the method beyond one city, we add four CVRPLIB instances~\cite{cvrplib}
(A-n32-k5, A-n33-k5, E-n101-k8, M-n200-k16) with their original demands,
capacities and Euclidean distances. For all seven instances, $\rho_{ij}$ and
therefore $t_{ij}$, $c_{ij}$ and $f_4$ are generated by the same procedure with
a fixed seed, so the objectives are defined identically across data sources.
The synthesized demands and road factors are a modelling assumption, not
measured data; we return to this in \cref{sec:conclusion}.

\begin{table}[t]
\centering
\caption{Benchmark instances. $n$: customers; $\lceil\sum q_i/Q\rceil$: lower
bound on the number of routes.}
\label{tab:instances}
\begin{tabular}{@{}llrrr@{}}
\toprule
Instance & Source & $n$ & $Q$ & $\lceil\sum q_i/Q\rceil$ \\
\midrule
A-n32-k5   & CVRPLIB & 31  & 100    & 5  \\
A-n33-k5   & CVRPLIB & 32  & 100    & 5  \\
E-n101-k8  & CVRPLIB & 100 & 200    & 8  \\
M-n200-k16 & CVRPLIB & 199 & 200    & 16 \\
route2\_199 & Sibiu  & 198 & 8\,469  & 6  \\
route3\_202 & Sibiu  & 201 & 8\,646  & 6  \\
route1\_334 & Sibiu  & 333 & 13\,988 & 6  \\
\bottomrule
\end{tabular}
\end{table}
